\documentclass[aspectratio=169]{beamer}

% Clean, professional Beamer setup
\usetheme{Madrid}
\usecolortheme{dolphin}
\setbeamertemplate{navigation symbols}{}
\setbeamertemplate{itemize item}{\color{structure.fg}$\blacktriangleright$}
\setbeamertemplate{itemize subitem}{\color{structure.fg}$\bullet$}

\usepackage[T1]{fontenc}
\usepackage[utf8]{inputenc}
\usepackage{graphicx}
\usepackage{booktabs}
\usepackage{array}

\title[Trade and Portfolio Analysis]{Trade and Portfolio Analysis Database System}
\subtitle{Database Design, Implementation, and Trade Analytics}
\author[Group 173 Yanggao]{Group 173 Yanggao\\Yingqi Tan \and Zhesheng Xu \and Hao Yang}
\institute[MISY330 Spring 2026]{MISY330 Database Design and Implementation\\Spring 2026}
\date{}

\newcommand{\queryimage}[1]{%
    \centering
    \includegraphics[width=\linewidth,height=0.36\textheight,keepaspectratio]{#1}
}

\newcommand{\largequeryimage}[1]{%
    \centering
    \makebox[\textwidth][c]{\includegraphics[width=1.08\textwidth,height=0.47\textheight,keepaspectratio]{#1}}
}

\newcommand{\mediumwidequeryimage}[1]{%
    \centering
    \makebox[\textwidth][c]{\includegraphics[width=0.92\textwidth,height=0.16\textheight,keepaspectratio]{#1}}
}

\newcommand{\widequeryimage}[1]{%
    \centering
    \makebox[\textwidth][c]{\includegraphics[width=1.08\textwidth,height=0.24\textheight,keepaspectratio]{#1}}
}

\newcommand{\prototypeimage}[1]{%
    \centering
    \includegraphics[width=\linewidth,height=0.64\textheight,keepaspectratio]{#1}
}

\newcommand{\algobox}[1]{%
    \begin{beamercolorbox}[rounded=true,sep=0.45em,wd=\linewidth]{block body}
    \textbf{Algorithm:} {\ttfamily\tiny #1}
    \end{beamercolorbox}
}

\begin{document}

% ------------------------------------------------
\begin{frame}
    \titlepage
\end{frame}

% ------------------------------------------------
\begin{frame}{Business Overview}
\begin{itemize}
    \item This project models a trade and portfolio analysis business that tracks stocks, exchanges, market data, strategies, factors, and trading records.
    \item A database is needed because trade analysis depends on structured relationships across companies, exchanges, portfolios, tax rules, and transaction history.
    \item The system supports practical tasks such as tax calculation, FIFO cost matching, factor-based stock screening, realized profit monitoring, and cross-market comparison.
    \item Business benefits include consistent record keeping, reusable query logic, faster reporting, and better decision support for portfolio management.
\end{itemize}
\end{frame}

% ------------------------------------------------
\begin{frame}{Database Scope and Main Entities}
\small
\begin{columns}[T,onlytextwidth]
    \column{0.48\textwidth}
    \begin{itemize}
        \item The database covers the full workflow from reference data to trade reporting.
        \item \texttt{Company}, \texttt{Exchange}, and \texttt{Stock} define the market structure.
        \item \texttt{MarketData} stores daily valuation and indicator information.
        \item \texttt{Portfolio} and \texttt{TradeRecord} capture holdings and executed transactions.
    \end{itemize}
    \column{0.52\textwidth}
    \begin{itemize}
        \item \texttt{Strategy} stores investment approaches used for trades.
        \item \texttt{Factor} stores screening dimensions such as PE ratio, dividend yield, RSI, and MACD.
        \item \texttt{FactorUsageRule} links factors to strategies through buy/sell rules.
        \item Together, these tables support tax analysis, FIFO matching, stock screening, and realized P\&L reporting.
    \end{itemize}
\end{columns}
\end{frame}

% ------------------------------------------------
\begin{frame}{How the System Works}
\small
\begin{itemize}
    \item The database is built in layers so each part of the business process has a clear role.
    \item \textbf{Reference layer:} \texttt{Company} stores issuer fundamentals, \texttt{Exchange} stores market rules and tax rates, and \texttt{Stock} connects each tradable ticker to its company and exchange.
    \item \textbf{Market layer:} \texttt{MarketData} stores date-based observations such as close price, EPS, dividend, RSI, and MACD for later valuation and factor analysis.
    \item \textbf{Decision layer:} \texttt{Factor}, \texttt{Strategy}, and \texttt{FactorUsageRule} model how screening rules are translated into buy and sell signals.
    \item \textbf{Transaction layer:} \texttt{Portfolio} identifies the investment account, and \texttt{TradeRecord} stores the executed trades with quantity, price, tax fields, selected strategy, and linked stock.
    \item \textbf{Analysis layer:} SQL queries combine these tables to calculate taxes, apply FIFO logic, summarize realized P\&L, and screen securities for investment decisions.
\end{itemize}
\end{frame}

% ------------------------------------------------
\begin{frame}{ER Diagram}
\begin{columns}[T,onlytextwidth]
    \column{0.38\textwidth}
    \scriptsize
    \begin{itemize}
        \item The ER diagram shows 9 entities, which satisfies the requirement of at least 8 entities.
        \item Primary keys identify each table, and foreign keys connect the business workflow from reference data to transactions.
        \item Major relationships include Exchange--Stock, Company--Stock, Stock--MarketData, Portfolio--TradeRecord, and Strategy--FactorUsageRule.
        \item \texttt{TradeRecord} is the central transaction table because each trade belongs to one portfolio, one strategy, and one stock.
    \end{itemize}
    \column{0.62\textwidth}
    \centering
    \includegraphics[width=\linewidth,height=0.70\textheight,keepaspectratio]{ER Diagram.png}
\end{columns}
\end{frame}

% ------------------------------------------------
\begin{frame}{Logical Data Model}
\begin{columns}[T,onlytextwidth]
    \column{0.38\textwidth}
    \small
    \begin{itemize}
        \item The logical model translates the ER design into implementation-ready tables and relationship directions.
        \item \texttt{FactorUsageRule} is the key associative table that resolves the many-to-many relationship between \texttt{Strategy} and \texttt{Factor}.
        \item \texttt{MarketData} uses a composite key of date and stock code to store time-series market observations.
        \item Optional tax fields in \texttt{TradeRecord} support later update queries for stamp duty and capital gain tax.
    \end{itemize}
    \column{0.62\textwidth}
    \centering
    \includegraphics[width=\linewidth,height=0.74\textheight,keepaspectratio]{Trade ER Model.png}
\end{columns}
\end{frame}

% ------------------------------------------------
\begin{frame}{Database Implementation}
\small
\begin{columns}[T,onlytextwidth]
    \column{0.55\textwidth}
    \begin{itemize}
        \item The schema was implemented in MySQL Workbench, and \texttt{GroupProject.sql} creates the database, tables, keys, and sample data.
        \item Parent tables are created first, then dependent tables such as \texttt{Stock}, \texttt{MarketData}, \texttt{FactorUsageRule}, and \texttt{TradeRecord}.
        \item The script loads sample records and keeps PK/FK links consistent so later tax, FIFO, and reporting queries can run directly on the same dataset.
    \end{itemize}
    \column{0.45\textwidth}
    \algobox{Create database\\
    Create parent tables and PK/FK rules\\
    Insert lookup data\\
    Load dependent tables\\
    Run update and analysis queries}
\end{columns}
\end{frame}

% ------------------------------------------------
\begin{frame}{Query 1: FIFO Cost Basis Check}
\scriptsize
\begin{itemize}
    \item \textbf{Idea:} match each sell trade to the earliest eligible buy in the same stock and portfolio.
    \item \textbf{Meaning:} FIFO cost basis is required for realized profit and capital gain tax.
    \item \textbf{Output:} NVDA sells match expected buy prices, while trade 20 has no prior buy match.
\end{itemize}
\algobox{Filter sell trades\\
Find earlier buy trades in same portfolio and stock\\
Return matched FIFO buy price for validation}
\largequeryimage{query 1.png}
\end{frame}

% ------------------------------------------------
\begin{frame}{Query 2: Updated Tax Results by Exchange}
\scriptsize
\begin{itemize}
    \item \textbf{Idea:} update taxes by linking each trade to exchange-specific rules.
    \item \textbf{Meaning:} profit reports are only reliable after the correct market tax logic is applied.
    \item \textbf{Output:} HKEX shows stamp duty, while NVDA and COST show capital gain tax on profitable US sells.
\end{itemize}
\algobox{Join TradeRecord + Stock + Exchange\\
Apply stamp duty rule by market and direction\\
Apply capital gain tax when realized gain > 0}
\largequeryimage{query 2.png}
\end{frame}

% ------------------------------------------------
\begin{frame}{Query 3: Total Realized P\&L by Portfolio and Currency}
\scriptsize
\begin{itemize}
    \item \textbf{Idea:} aggregate realized profit for sell trades after tax.
    \item \textbf{Meaning:} management needs a quick portfolio-level performance summary by currency.
    \item \textbf{Output:} \texttt{Main\_Aggressive} shows total realized P\&L of 2067.75 in USD.
\end{itemize}
\algobox{Use sell trades only\\
Compute net realized profit after tax\\
Group by portfolio and trade currency}
\mediumwidequeryimage{query 3.png}
\end{frame}

% ------------------------------------------------
\begin{frame}{Query 4: Transaction-Level Realized Profit Detail}
\scriptsize
\begin{itemize}
    \item \textbf{Idea:} show realized profit for each sell trade inside one portfolio.
    \item \textbf{Meaning:} managers need trade-level audit detail behind the summary report.
    \item \textbf{Output:} the portfolio contains one NVDA gain, one AAPL loss, and one break-even NVDA exit.
\end{itemize}
\algobox{Select one portfolio\\
Match sell price with FIFO buy price and taxes\\
Return net realized profit by transaction}
\widequeryimage{query 4.png}
\end{frame}

% ------------------------------------------------
\begin{frame}{Query 5: Negative-Growth Companies with PE Ratio}
\scriptsize
\begin{itemize}
    \item \textbf{Idea:} screen stocks whose companies have negative profit growth and compute PE ratio.
    \item \textbf{Meaning:} this helps identify weak-growth firms and compare valuation across markets.
    \item \textbf{Output:} the result includes Chagee, BYD, JPMorgan Chase, and China Mobile.
\end{itemize}
\algobox{Join Company + Stock + Exchange + MarketData\\
Keep rows with negative profit growth\\
Calculate PE ratio from close price and EPS}
\widequeryimage{query 5.png}
\end{frame}

% ------------------------------------------------
\begin{frame}{Query 6: Portfolios with Repeated Sell Activity}
\scriptsize
\begin{itemize}
    \item \textbf{Idea:} find portfolios with at least two sell trades in the same currency.
    \item \textbf{Meaning:} this identifies active accounts and summarizes tax-adjusted sale proceeds.
    \item \textbf{Output:} three portfolio-currency groups qualify, including \texttt{Main\_Aggressive} in USD.
\end{itemize}
\algobox{Use sell trades only\\
Group by portfolio and currency\\
Keep groups with COUNT(*) >= 2}
\widequeryimage{query 6.png}
\end{frame}

\begin{frame}{Technical Highlights}
\begin{itemize}
    \item Primary-key and foreign-key structure organizes the schema around clean links between portfolios, trades, stocks, exchanges, strategies, factors, and market data.
    \item Multi-table joins support realistic business reporting instead of isolated table lookups.
    \item Nested subqueries and a CTE with cumulative logic are used to model FIFO cost basis for sell transactions.
    \item Query design is business-oriented: taxes are exchange specific, realized profit is portfolio specific, and screening logic uses both fundamentals and market metrics.
    \item The overall logic moves from raw records to decision-ready outputs: base tables store data, update queries calculate tax fields, nested queries resolve FIFO matching, and summary queries transform those results into reports.
    \item The project reflects real trade-analysis ideas such as stamp duty, capital gain tax, FIFO matching, valuation ratios, and realized profit measurement.
\end{itemize}
\end{frame}

% ------------------------------------------------
\begin{frame}{Bonus Prototype Demo}
\begin{itemize}
    \item We completed the bonus task by building a simple local Flask application connected to the MySQL trade database.
    \item The implementation uses:
    \begin{itemize}
        \item Flask for routing and page rendering
        \item PyMySQL for direct database access
        \item HTML templates for dashboard, company browser, trade CRUD, reports, and connection settings
        \item SQL queries that reuse the project tables instead of mock data
    \end{itemize}
    \item The prototype is local only, with no authentication and no deployment, which matches the project bonus description.
    \item This extension shows that the database is not only well designed in SQL, but also usable in a small application workflow.
\end{itemize}
\end{frame}

% ------------------------------------------------
\begin{frame}{Prototype Implementation Logic}
\small
\begin{itemize}
    \item The prototype follows the same logic as the database project and turns the schema into a small usable application.
    \item Step 1: start the Flask app locally in PowerShell with \texttt{python app.py}.
    \item Step 2: open the connection settings page and enter host, port, username, password, target schema, and SSL option.
    \item Step 3: after the connection test succeeds, the app runs SQL against the actual MySQL database and displays live results.
    \item Step 4: the interface supports browsing reference data, viewing summarized dashboard metrics, reading query-based reports, and editing \texttt{TradeRecord} rows.
    \item This design keeps the architecture simple: browser $\rightarrow$ Flask routes $\rightarrow$ PyMySQL queries $\rightarrow$ MySQL database.
\end{itemize}
\end{frame}

% ------------------------------------------------
\begin{frame}{Prototype Setup and Secure Connection}
\begin{columns}[T,onlytextwidth]
    \column{0.48\textwidth}
    \scriptsize
    \textbf{Local startup}
    \begin{itemize}
        \item The Flask server is launched from PowerShell inside the project folder.
        \item The app runs locally at \texttt{http://127.0.0.1:5000}.
        \item This keeps the demo simple and classroom ready.
    \end{itemize}
    \includegraphics[width=\linewidth,height=0.44\textheight,keepaspectratio]{powershell setting.png}
    \column{0.52\textwidth}
    \scriptsize
    \textbf{Database connection}
    \begin{itemize}
        \item A settings page lets us configure the real MySQL target without editing code.
        \item The connection uses the course server, project schema, user credentials, and SSL/TLS because Workbench confirmed secure access.
        \item This page also helped us debug connection issues step by step during implementation.
    \end{itemize}
    \includegraphics[width=\linewidth,height=0.44\textheight,keepaspectratio]{website0.png}
\end{columns}
\end{frame}

% ------------------------------------------------
\begin{frame}{Prototype Running Result}
\begin{columns}[T,onlytextwidth]
    \column{0.5\textwidth}
    \scriptsize
    \textbf{Connection success and dashboard}
    \begin{itemize}
        \item After a successful test, the app confirms the database target and SSL status.
        \item The dashboard immediately shows live counts for companies, stocks, portfolios, and trades.
        \item This proves the Flask app is reading real data from the database.
    \end{itemize}
    \includegraphics[width=\linewidth,height=0.46\textheight,keepaspectratio]{website1.png}
    \column{0.5\textwidth}
    \scriptsize
    \textbf{Data display and reporting}
    \begin{itemize}
        \item The dashboard continues with recent trades and a sell-amount-after-tax summary.
        \item Other pages also support company browsing, trade CRUD, and query-based reports.
        \item Together, the prototype demonstrates a complete local front-end workflow on top of the project database.
    \end{itemize}
    \includegraphics[width=\linewidth,height=0.46\textheight,keepaspectratio]{website2.png}
\end{columns}
\end{frame}

% ------------------------------------------------
\begin{frame}{Challenges and Improvements}
\begin{columns}[T,onlytextwidth]
    \column{0.5\textwidth}
    \textbf{Challenges}
    \begin{itemize}
        \item Designing a schema that connects market data, strategy rules, and trade records without redundancy
        \item Implementing FIFO-related logic with SQL subqueries and trade ordering
        \item Handling exchange-specific tax rules across HK, mainland China, and US markets
        \item Maintaining consistent key values across all sample records
    \end{itemize}
    \column{0.5\textwidth}
    \textbf{Future Improvements}
    \begin{itemize}
        \item Add views for reusable reporting outputs
        \item Add stored procedures or functions for tax and FIFO calculation
        \item Extend the Flask prototype with more pages, filters, and form validation
        \item Add authentication and user-role control if the prototype is expanded beyond a classroom demo
        \item Extend market data history for richer performance analytics
    \end{itemize}
\end{columns}
\vspace{0.4em}
\small Current version already includes a working local Flask prototype, but it is still intentionally lightweight and focused on classroom demonstration.
\end{frame}

% ------------------------------------------------
\begin{frame}{Conclusion}
\begin{itemize}
    \item This database integrates business reference data, market data, strategy rules, and transaction records into one consistent trade-analysis system.
    \item The project supports meaningful business tasks such as tax calculation, FIFO cost matching, stock screening, and portfolio profit reporting.
    \item Overall, the system demonstrates how database design and SQL implementation can support practical portfolio trading analysis.
\end{itemize}
\end{frame}

% ------------------------------------------------
\begin{frame}{Contribution and AI Usage}
\small
\begin{itemize}
    \item \textbf{Yingqi Tan:} business background, presentation organization, and final review.
    \item \textbf{Zhesheng Xu:} schema design, ER/logical modeling, and MySQL implementation.
    \item \textbf{Hao Yang:} SQL query design, result verification, and slide integration.
    \item \textbf{Shared contribution:} all members discussed the database logic, reviewed the project outputs, and understood the complete workflow.
    \item \textbf{AI usage statement:} AI tools were used appropriately for debugging, wording refinement, and presentation drafting support; the database design, SQL logic, and final verification were completed and checked by the group.
\end{itemize}
\end{frame}

% ------------------------------------------------
\begin{frame}
\centering
\vspace{2.2cm}
{\Huge Any Questions?}
\vspace{0.8cm}

{\large Thank you for listening.}
\end{frame}

\end{document}
